% Generated by Naberius. Content is escaped; shell escape is disabled.
\documentclass[12pt,a4paper,oneside,openany]{report}
\usepackage{fontspec}
\defaultfontfeatures{Ligatures=TeX}
\setmainfont[Path=fonts/,BoldFont=lmroman10-bold.otf,ItalicFont=lmroman10-italic.otf,BoldItalicFont=lmroman10-bolditalic.otf]{lmroman10-regular.otf}
\usepackage[spanish]{babel}
\usepackage[margin=25.0mm,headheight=16pt]{geometry}
\usepackage{graphicx,amsmath,amssymb,tabularx,booktabs}
\usepackage{fancyhdr}
\usepackage[unicode,hidelinks]{hyperref}
\hypersetup{pdftitle={El oficio de escribir},pdfauthor={Documento de demostración},pdfkeywords={}}
\addto\captionsspanish{\renewcommand{\figurename}{Figura}\renewcommand{\tablename}{Tabla}\renewcommand{\contentsname}{Índice}}
\setcounter{secnumdepth}{3}
\counterwithout{figure}{chapter}
\counterwithout{table}{chapter}
\counterwithout{equation}{chapter}
\pagestyle{fancy}
\fancyhf{}\fancyhead[L]{\small Prueba de \LaTeX{}}\fancyfoot[C]{\small Naberius\quad \thepage}\renewcommand{\headrulewidth}{0.4pt}
\fancypagestyle{plain}{\fancyhf{}\fancyhead[L]{\small Prueba de \LaTeX{}}\fancyfoot[C]{\small Naberius\quad \thepage}\renewcommand{\headrulewidth}{0.4pt}}
\linespread{1.4}
\setlength{\parindent}{0.0mm}
\setlength{\parskip}{7.0pt}
\begin{document}
\chapter*{Resumen}
Prueba tipográfica de \LaTeX{}: español, áéíóú y referencias.
\fontsize{12.0}{14.399999999999999}\selectfont
{\LARGE\bfseries El oficio de escribir\par}

{\large Del primer borrador a una idea que se entiende.\par}

\chapter{Una idea necesita espacio}\label{bAD694EE6-5445-4C8C-B4C6-60B27B440AE0}

Escribir es dar forma a lo que pensamos. Este documento de demostración muestra cómo organizar un texto, enlazar sus fuentes y cuidar su presentación sin interrumpir la escritura.

\section{Estructura antes que apariencia}\label{bC99A71E5-0CCD-412B-84F4-F0D0838E62E3}

Usa los estilos de la barra para expresar qué significa cada párrafo. El esquema se construye a partir de tus títulos; la plantilla se ocupa de la tipografía.

Inserta una tabla, una figura o una ecuación desde el menú Insertar.

\begin{itemize}
\item Guarda tu trabajo como .nab: contenido, imágenes y referencias viajan juntos.

\end{itemize}
\chapter{Un documento, dos maneras de verlo}\label{b4B3EFB2D-E8F5-498A-A0C0-59ABB5014B69}

Alterna entre páginas y lienzo continuo. El texto y el historial son los mismos; la paginación definitiva corresponde al PDF compuesto con \LaTeX{}.

\LaTeX{} y \TeX{} conservan su composición tipográfica.

\appendix
\chapter{Datos complementarios}\label{b941AEF8A-FFA0-4FD7-AC54-F2CC1871DD6A}

\par\medskip\noindent\rule{\linewidth}{0.4pt}\par\medskip

\begin{table}[htbp]\centering\caption{Resultados de prueba}\label{n10A1D6DE-2589-4B7F-8568-CD6137F7206E}
\begin{tabularx}{\linewidth}{|>{\raggedright\arraybackslash}X|>{\raggedright\arraybackslash}X|}\hline
\textbf{Variable} & \textbf{Valor} \\ \hline
 &  \\ \hline
\end{tabularx}\par{\small }\end{table}

\begin{equation}\sqrt{\alpha  + \beta }\label{n75EE53E5-393E-45D9-A3FA-5B7779C41A80}\end{equation}

\end{document}